\chapter{Literature Review}
\label{ch:literature}

The literature review tells the story of the work your thesis builds on. Group studies by
theme, compare them, and show the gap that your work fills.

\section{Citing Sources with Page Numbers}
UKH requires a page number in every in-text citation, whatever the referencing style
(Guideline 1c). Put the page in square brackets before the key:

\begin{verbatim}
\parencite[p.~9461]{lewis2020rag}
\textcite[p.~425]{jarvelin2002cumulated}
\parencite[pp.~758--759]{cormack2009reciprocal}
\end{verbatim}

These give \parencite[p.~9461]{lewis2020rag}, then \textcite[p.~425]{jarvelin2002cumulated},
and \parencite[pp.~758--759]{cormack2009reciprocal}.

Direct quotations longer than about forty words go in a \texttt{quote} environment, which
is single spaced as the guidelines require:

\begin{quote}
An indented quotation is set with single line spacing, as Guideline 3d asks. Give the
source and page number at the end \parencite[p.~2]{mahmudi2023meter}.
\end{quote}

\section{Writing Kurdish Text}
Sorani Kurdish in Arabic script works out of the box. Inside an English sentence use
\verb|\textkurdish{...}|, for example \textkurdish{زانکۆی کوردستان هەولێر}
(University of Kurdistan Hewlêr). For a whole right-to-left paragraph, such as a quoted
line of poetry, use the \texttt{kurdish} environment:

\begin{kurdish}
ئەمە نموونەیەکە بۆ نووسینی دەقی کوردی لە ناو نامەی ماستەردا.
پیتە تایبەتەکانی کوردی: ڕ ڵ ێ ۆ ە ڤ.
\end{kurdish}

Kurdish studies relevant to this kind of work include information retrieval
\parencite[p.~3]{esmaili2014kurdish} and poetry analysis \parencite[p.~1]{mahmudi2023meter}.

\section{Summary}
Summarise the gap in the literature that your thesis addresses.
